% Section 4: R x C under LCE (PHASE2a)
\section{R$\times$C under linear contextual effects}\label{sec:rxc}

Source: PHASE2a (Appendix~\ref{app:phase2a}), which builds on PHASE1
(Appendix~\ref{app:phase1}). Every proof in PHASE2a is written in numbered
steps with a justification tag; the tag legend is PHASE2a \S0.3 and the citation
check is \S0.4. Numerical and symbolic checks: \texttt{verify\_phase2a.py},
32/32 PASS. The model $\mathcal M(H;A)$ is the set of laws of $(W,\beta)$ with
(S1) $\beta\in\Bset$ and $T=\beta^\top X$; (S2) $\E[\beta\mid V]=M_H(X)$ on
$\{X\in A\}$; (S3) $M_H(x)\in\Bset$ for every $x\in A$.

\subsection{The unit polytope and what the data identify}

\begin{result}[Lemma 0 --- PHASE2a \S1]\label{res:p2a-l0}
\status{Proved.} For all $x\in\Delta^{R-1}$, $t\in\Delta^{C-1}$: (a)
$S(x,t)\ne\emptyset$; (b) it is compact and convex; (c) $(x,t)\mapsto S(x,t)$ has
closed graph, is upper hemicontinuous and Borel; (d) $h_S$ is sublinear and
$\sqrt R$-Lipschitz in $u$, concave in $t$, Borel; (e) $S$ is the intersection
of its supporting half-spaces.
\end{result}
\begin{plain}
Each unit's set of possible rate tables is never empty (``everyone in the unit
behaves like the unit average'' always fits), and it is a well-behaved convex
polytope. These are the regularity facts every later proof uses.
\end{plain}
\proofref{Appendix~\ref{app:phase2a}, \S1, steps 1--13}

\begin{result}[Lemma 1, identification decomposition --- PHASE2a \S2]\label{res:p2a-l1}
\status{Proved.} Under the model and A5: (a) $\E[T_c\mid X=x]=x^\top Q_cx$ and
$Q_c$ is unique; (b) $\sum_cH_c=J$, hence $\sum_cQ_c=J$ and $\sum_cK_c=0$;
(c) every $K'\in\Kspace$ gives the same regression, and
$\dim\Kspace=(C-1)R(R-1)/2$.
\end{result}
\begin{plain}
The data reveal the symmetric part $Q$ of the contextual-effect matrices, and
nothing about the skew part $K$. For each pair of groups and each outcome, the
data identify the sum ``effect of $s$-share on group $r$ plus effect of $r$-share
on group $s$'' but not how it splits. A $2\times2$ table has one unknown, a
$3\times3$ table six.
\end{plain}
\begin{sketch}
Take conditional expectations of the accounting identity: $\E[T_c\mid X]=X^\top
H_cX$. A skew matrix vanishes in a quadratic form, so only $Q_c$ appears. A5
makes the quadratic form determine $Q_c$. Row sums of $\beta$ equal one, which
under A5 forces $\sum_cH_c=J$.
\end{sketch}
\proofref{Appendix~\ref{app:phase2a}, \S2, (a) steps 1--8, (b) steps 1--7, (c) steps 1--3}
\begin{revnote}
(R-01, R-02, P-B39.) The derivation of $\sum_cH_c=J$ now uses A5 rather than an
unassumed open set. A5 is a matrix rank condition stated on
$\supp_AP_X$, the support of $P_X$ restricted to $A$; the earlier
$\supp P_X\cap A$ is wrong when $A$ carries less than full mass
($X\sim U[0,1]$, $A=\{1/2\}$).
\end{revnote}

\paragraph{Remarks 1.0--1.2 (PHASE2a \S2).} \emph{Remark~1.0}: A2 cannot be
validated by goodness of fit; Prop.~9 exhibits nonlinear departures that leave
the data unchanged (moved forward after review, R-16). \emph{Remark~1.1}
(proved): for $R=2$, A5 holds iff there are at least three distinct
compositions; for $R=3$ six points off a common conic are needed, and five are
never enough. \emph{Remark~1.2} (proved): column-by-column OLS on the quadratic
monomials gives $\sum_c\hat Q_c=J$ exactly, so no constrained regression is
needed.

\begin{result}[Lemma 2, the estimand --- PHASE2a \S3]\label{res:p2a-l2}
\status{Proved.} Under (S1)--(S2) with $P(X\in A)=1$,
\[
B^*_{rc}=\frac1{\mu_r}\Big[\sum_sQ_{c,rs}\mu_{rs}+\sum_{s\ne r}K_{c,rs}\mu_{rs}\Big],
\]
affine in $K$, with coefficient $\mu_{rs}/\mu_r\ge0$ on the ambient entry $K_{c,rs}$.
\end{result}
\begin{plain}
The target is a known number plus a weighted sum of the unknown $K$ entries. The
weights $\mu_{rs}/\mu_r$ measure how much group $r$ lives alongside group $s$:
the more mixed the units, the more the unknown split matters.
\end{plain}
\proofref{Appendix~\ref{app:phase2a}, \S3, steps 1--7; $2\times2$ check against JKST (16)}
\begin{revnote}
(R-03, R-05.) The coefficient is stated in ambient coordinates, with the free
coordinates in Prop.~7. $P(X\in A)=1$ is now explicit; Remark~2.1 and
Prop.~7$'$ treat units outside $A$.
\end{revnote}

\subsection{The sharp identified set}

\begin{result}[Theorem 1, characterisation of $\ThI$ --- PHASE2a \S4]\label{res:p2a-t1}
\status{Proved} (measure-theoretic steps rest on the cited theorems of PHASE2a
\S0.4). Assume A1, A5, standing integrability, $H=Q(P)+K$ with $K\in\Kspace$.
Then $K\in\ThI(P;A)$ iff (S3) holds and, outside one $P_V$-null set, for every
$v=(x,n)$ with $x\in A$,
\[
\langle u,M_H(x)\rangle\le\psi(u;v):=\int h_S(u;x,t)\,P_{T\mid V=v}(dt)\quad\text{for all }u\in\mathbb R^{R\times C}. \tag{T1}
\]
$\ThI$ is convex and closed; it is compact when $\operatorname{affdim}(A)=R-1$,
which A5 implies.
\end{result}
\begin{plain}
A value of the unknown $K$ is consistent with the data exactly when, at every
composition and size, the implied average rate table could be the average of
rate tables that each fit some unit's margins. It is the strongest possible
use of the data under the model. It needs the whole conditional distribution of
outcomes given composition and size, which in real data must be smoothed, so it
is not what the package computes.
\end{plain}
\begin{sketch}
Necessity: each unit's $\beta$ lies in its polytope, so its conditional mean lies
in the conditional selection expectation; one common null set for all $u$ is
obtained from a countable dense set and Lipschitz continuity. Sufficiency: the
set of measurable selections is weakly compact in $L^1$ (Dunford--Pettis,
Mazur); its image under $\E[\cdot\mid V]$ is weakly compact and convex; if the
target mean were outside it, Hahn--Banach would separate it, contradicting
(T1) evaluated at the separating direction (a measurable argmax attains the
support function). No auxiliary randomisation is needed (Remark~4.1). The random-set formulation follows Beresteanu, Molchanov and Molinari~\citep{bmm2011}.
\end{sketch}
\proofref{Appendix~\ref{app:phase2a}, \S4: (i)$\Rightarrow$(ii) steps 1--11; facts F1--F7; (ii)$\Rightarrow$(i) steps 1--22; convexity and compactness steps 1--4}
\begin{revnote}
(R-07, R-08, R-09.) Compactness needs $\operatorname{affdim}(A)=R-1$, not
``nonempty relative interior''. The sufficiency proof was expanded into
facts F1--F7 with checked citations, and a paragraph explains why (T1) must hold
for all $u$ outside one null set.
\end{revnote}

\paragraph{Remark 4.2 (continuous versus integer rates; PHASE2a \S4).}
Sharpness is relative to real-valued latent rates. Part~(1), that every vertex of
the transportation polytope is an integer table, is a standard fact checked on
200 cases. Part~(2), that randomising over integer vertices would leave the
sharp set unchanged, is \emph{only a sketch} (not verified) and is not used
anywhere (R-12).

\begin{result}[Corollary 1, sharp bounds for a linear target --- PHASE2a \S4]
\status{Proved.} If $P(X\in A)=1$ and $\ThI\ne\emptyset$, the sharp set for
$\varphi$ is $[\min_{\ThI}\varphi(K),\max_{\ThI}\varphi(K)]$.
\end{result}
\begin{plain}
Because the unknown enters the target linearly and the consistent set is convex
and compact, the answer is always an interval.
\end{plain}
\proofref{Appendix~\ref{app:phase2a}, \S4, Corollary~1 steps 1--3}

\subsection{The outer set: reduction, nesting, sharpness for the regression}

\begin{result}[Lemma 3 --- PHASE2a \S5]\label{res:p2a-l3}
\status{Proved.} (a) For a polytope $A$, linear constraints on $A$ reduce to its
vertices. (b) If $\operatorname{affdim}(A)=R-1$, $\ThO$ is bounded (a polytope);
otherwise it is a closed convex polyhedron. (c) Under (S1)--(S2), if
$A\subseteq\conv(\supp_AP_X)$, (S3) holds automatically, in particular for the
population hull.
\end{result}
\begin{plain}
Only the most extreme units matter for the constraints. If the region $A$ is
the range of the data, probability validity of the straight lines is automatic;
if $A$ extrapolates beyond the data, validity there is an extra assumption.
\end{plain}
\proofref{Appendix~\ref{app:phase2a}, \S5, steps 1--9; Remarks~5.1--5.2 give the counterexamples}
\begin{revnote}
(R-04, R-07, P-B39.) (c) was restated: (S3) is automatic only when
$A\subseteq\conv(\supp_AP_X)$. Remark~5.1 (a singleton inside the hull with no
mass) and Remark~5.2 (a singleton $A$ makes $\ThO$ unbounded; A5 rules it out)
were added. Checks in \texttt{verify\_phase2a.py} fail on the earlier statements.
\end{revnote}

\begin{result}[Proposition 1, nesting --- PHASE2a \S5]\label{res:p2a-p1}
\status{Proved.} $\ThI(P;A)\subseteq\ThX(P;A)\subseteq\ThO(P;A)$.
\end{result}
\begin{plain}
Using more of the data can only narrow the set. The outer set uses only the
regression, so it is the widest of the three.
\end{plain}
\begin{sketch}
Conditioning on $(X,N)$ implies conditioning on $X$ (tower property). Conditional
expectations of valid rate tables are valid rate tables, which is JKST's Jensen
step.
\end{sketch}
\proofref{Appendix~\ref{app:phase2a}, \S5, Proposition~1 steps 1--5}

\begin{result}[Proposition 2, $\ThO$ is sharp for the regression function --- PHASE2a \S5]\label{res:p2a-p2}
\status{Proved.} Let $\mathcal P(P)$ be the laws with the same law of $(N,X)$
and the same regression $\E[T\mid X]$. If $P(X\in A)=1$,
$\ThO(P;A)=\bigcup_{P'\in\mathcal P(P)}\ThI(P';A)$.
\end{result}
\begin{plain}
The outer bound is exactly what can be learned from the regression alone. It is
the sharp set for any data set in which units with the same composition have
identical outcomes. Calling it ``sharp'' without that qualifier is wrong.
\end{plain}
\begin{sketch}
Given $K\in\ThO$, set every unit's rates equal to the implied mean $M_H(X)$;
this is a valid law with the same regression, and $K$ is in its sharp set.
\end{sketch}
\proofref{Appendix~\ref{app:phase2a}, \S5, Proposition~2 steps 1--10}

\begin{result}[Corollary 2 and Proposition 3 --- PHASE2a \S5]
\status{Proved.} If $T=\E[T\mid X]$ a.s., then $\ThI=\ThO$ (Cor.~2). More
generally (Prop.~3, a \emph{sufficient} condition), if $\E[T\mid V]=\E[T\mid X]$
and, for a.e.\ $v$ and every $u$, $t\mapsto h_S(u;x,t)$ is affine on the
convex hull of the conditional support of $T$, then $\ThI=\ThO$.
\end{result}
\begin{plain}
The regression-only bound loses nothing unless units with similar composition
are so different that some sit at a Duncan--Davis corner and others do not
(Remark~3.1, the $2\times2$ reading). In 400 random valid $2\times2$ designs the
inclusion was strict in 46.
\end{plain}
\proofref{Appendix~\ref{app:phase2a}, \S5, Corollary~2 steps 1--3; Proposition~3 steps 1--7}

\subsection{Closed forms and the linear program}

\begin{result}[Proposition 4, equivalence with JKST --- PHASE2a \S6.1]\label{res:p2a-p4}
\status{Proved.} In $2\times2$ with $A=[l,u]$: (a) $Q_{1,22}=w_0$,
$Q_{1,12}=w_0+c_1/2$, $Q_{1,11}=w_0+c_1+d_1$, $K=c_1/2-w_1$; (b) $\ThO$ is JKST
Prop.~2; (c) with $A=\Delta^1$ it is JKST Prop.~3; (d) with $P(l\le X\le u)=1$,
$B^*(K)$ is JKST eq.~(16).
\end{result}
\begin{plain}
The R$\times$C machinery contains the published $2\times2$ method exactly; the
package's existing bounds are a special case.
\end{plain}
\proofref{Appendix~\ref{app:phase2a}, \S6.1, derivations (a)--(d); numerically to $10^{-16}$}

\begin{result}[Proposition 5, $A$ the whole simplex --- PHASE2a \S6.2]
\status{Proved.} With $A=\Delta^{R-1}$, $\ThO=\{|K_{c,rs}|\le Q_{c,rs}\}$; it is
non-empty iff every $Q_{c,rs}\ge0$; each $K_{c,rs}$ ranges over
$[-m_{c,rs},m_{c,rs}]$ with $m_{c,rs}=\min\{Q_{c,rs},1-Q_{c,rs}\}$, and $B^*$
has a closed form.
\end{result}
\begin{plain}
If the straight lines are trusted all the way to pure units, two hypothetical
probabilities with known average $Q$ can split it by at most
$\pm\min(Q,1-Q)$. This is narrower than the hull-based bound but assumes
extrapolation.
\end{plain}
\proofref{Appendix~\ref{app:phase2a}, \S6.2, steps 1--8; checked against an LP in 200 tables}

\begin{result}[Lemma 4, box $\cap$ hyperplane --- PHASE2a \S6.3]
\status{Proved.} For $\mathcal P=\{\ell\le\kappa\le h,\sum_c\kappa_c=0\}$:
(a) non-empty iff $\sum\ell\le0\le\sum h$; (b) coordinate projections in closed
form; (c) linear objectives are maximised by a greedy rule.
\end{result}
\proofref{Appendix~\ref{app:phase2a}, \S6.3, steps 1--6}

\begin{result}[Proposition 6, two rows, any number of columns --- PHASE2a \S6.4]
\status{Proved.} For $R=2$ and $A=[l,u]$ with $0<l<u<1$, $\ThO=\{\ell\le\kappa\le
h,\ \sum_c\kappa_c=0\}$ with closed-form $\ell_c,h_c$; target bounds follow by
Lemma~4.
\end{result}
\proofref{Appendix~\ref{app:phase2a}, \S6.4, steps 1--6; checked against an LP in 300 cases}

\begin{result}[Corollary 6.1, joint versus collapsed --- PHASE2a \S6.4]
\status{Proved.} For $R=2$: (a) the joint interval for $\kappa_c$ lies inside the
collapsed one; (b) they are equal when, for each row, the $Q_{c',rr}$ ($c'\ne c$)
share a sign; (c) strict tightening therefore requires some negative
extrapolated pure-poll rate. For $R\ge3$ the joint bound can be strictly tighter
with all rates valid; a characterisation is open.
\end{result}
\begin{plain}
With two groups, fitting all outcome columns jointly helps only if the straight
lines would give an impossible rate at some pure unit. With three or more groups,
joint fitting can help even when everything is valid.
\end{plain}
\proofref{Appendix~\ref{app:phase2a}, \S6.4, Corollary~6.1 steps 1--7}
\begin{revnote}
(R-14; PHASE2a \S10 item~1.) PHASE1's general claim of strict joint gain was
withdrawn; the exact sign condition replaced ``strong contextual effects''.
\end{revnote}

\begin{result}[Proposition 7, the general LP --- PHASE2a \S6.5]\label{res:p2a-p7}
\status{Proved.} With free coordinates $k_{c,rs}$ ($c<C$, $r<s$),
$\ThO=\{k:Gk\le g\}$, one row per (vertex, row, column), and
$\varphi(k)=\varphi_0+\gamma^\top k$ with explicit $\varphi_0$ and
$\gamma_{c,rs}=\mu_{rs}[(a_{rc}-a_{rC})/\mu_r-(a_{sc}-a_{sC})/\mu_s]$. Imposing
the rows at every observed $X_i$ equals imposing them at the hull vertices.
\end{result}
\begin{plain}
The computable bound is a small linear program: one variable per unknown
direction, one constraint per extreme unit, group and outcome. This is the
\emph{outer} set, not an LP form of the sharp set.
\end{plain}
\proofref{Appendix~\ref{app:phase2a}, \S6.5, steps 1--5}
\begin{revnote}
(R-06.) The closing paragraph on regularity was wrong: the value function is
concave and piecewise linear in $Q$ for fixed $\mu$, convex in $\gamma$ for
fixed $Q$, and in general neither jointly ($V=\mu Q$). Lemma~P3.1 is the
fixed-design case. Check: \texttt{verify\_phase2a.py} (midpoint test).
\end{revnote}

\begin{result}[Proposition 7$'$, units outside $A$ --- PHASE2a \S6.6]
\status{Proved.} If $P(X\in A)>0$ only, the bound is the LP over units in $A$
plus the support-function (Duncan--Davis) terms of the units outside $A$; each
endpoint is attained. This is what \texttt{range = "trim"} computes for single
cells.
\end{result}
\begin{plain}
Units outside the region where straight lines are trusted contribute their
model-free range; nothing links them to the model.
\end{plain}
\proofref{Appendix~\ref{app:phase2a}, \S6.6, steps 1--6}
\begin{revnote}
(R-05; decision D2.) New result written after review to document the existing
\texttt{trim} option.
\end{revnote}

\subsection{Refutability and undetectable violations}

\begin{result}[Proposition 8, non-emptiness --- PHASE2a \S7]
\status{Proved.} If $P$ is generated by the model with $P(X\in A)=1$, then the
true $K_0\in\ThI\subseteq\ThO$ and $B^*_{rc}\in[\varphi_L,\varphi_U]\cap DD^*_{rc}$.
An empty $\ThO$, or an empty intersection with $DD^*$, refutes A1--A5 jointly.
\end{result}
\begin{plain}
If the model is right, the population bound is never empty. An empty bound is
evidence against the model (in samples it can also be noise; Phase~3 turns it
into a test).
\end{plain}
\proofref{Appendix~\ref{app:phase2a}, \S7, steps 1--5}

\begin{result}[Proposition 9, what the data can never rule out --- PHASE2a \S7]\label{res:p2a-p9}
\status{Proved.} If $\Delta(x)$ satisfies $x^\top\Delta(x)=0$, $\Delta(x)\mathbf
1=0$ and keeps $M_H+\Delta$ valid on the support, then $\beta=M_H(X)+\Delta(X)$
and $\beta'=M_H(X)$ give the same $P$ but different $B^*$. Under A5 the linear
such $\Delta$ are exactly $M_{K'}$, $K'\in\Kspace$.
\end{result}
\begin{plain}
There are nonlinear contextual effects that cancel in the aggregate and leave
the data identical, yet change the transition rates. The linearity assumption is
therefore a genuine leap that no fit statistic can confirm. Phase~3 bounds its
consequences with the sensitivity parameter $\eta$ (Prop.~S1).
\end{plain}
\proofref{Appendix~\ref{app:phase2a}, \S7, Proposition~9 steps 1--4}

\paragraph{Illustration (PHASE2a \S8).} In a noise-free BC $3\times3$ design with
2000 polls, the $\ThO$ bounds on loyalty, defection, demobilisation and
mobilisation have width ratios 0.52--0.64 relative to DD and all contain the
truth. They remain wide because six numbers are unidentified.

\paragraph{Relation to Cross--Manski (PHASE2a \S12).} Cross and
Manski~\citep{crossmanski2002} treat Goodman-type constancy (an instrument) and
obtain a point; LCE is weaker, and the R$\times$C bounds sit between their
data-alone region and their point.
